\documentclass[12pt, letterpaper]{article}

% ============================================================
%  Paquetes
% ============================================================
\usepackage[utf8]{inputenc}
\usepackage[T1]{fontenc}
\usepackage[spanish,es-noshorthands]{babel}
\usepackage[letterpaper,margin=2.5cm,headheight=15pt]{geometry}
\usepackage{graphicx}
\usepackage{xcolor}
\usepackage{booktabs}
\usepackage{longtable}
\usepackage{array}
\usepackage{enumitem}
\usepackage{caption}
\usepackage{fancyhdr}
\usepackage{titlesec}
\usepackage{listings}
\usepackage{amsmath}
\usepackage{amssymb}
\usepackage[hidelinks,colorlinks=true,linkcolor=winorange,urlcolor=winorange,citecolor=winorange]{hyperref}

% ============================================================
%  Estilo
% ============================================================
\definecolor{winorange}{HTML}{FF6B00}
\definecolor{windark}{HTML}{1A1A1A}
\definecolor{codebg}{HTML}{F5F5F5}
\definecolor{codegray}{HTML}{6A737D}
\definecolor{codekw}{HTML}{C7361B}
\definecolor{codestr}{HTML}{1A7F37}

\titleformat{\section}{\Large\bfseries\color{windark}}{\thesection}{1em}{}
\titleformat{\subsection}{\large\bfseries\color{winorange}}{\thesubsection}{1em}{}

\pagestyle{fancy}
\fancyhf{}
\lhead{\textcolor{windark}{WinSports}}
\rhead{\textcolor{codegray}{Informe Final}}
\cfoot{\thepage}
\renewcommand{\headrulewidth}{0.4pt}

\lstdefinestyle{winstyle}{
  backgroundcolor=\color{codebg},
  basicstyle=\ttfamily\footnotesize,
  keywordstyle=\color{codekw}\bfseries,
  stringstyle=\color{codestr},
  commentstyle=\color{codegray}\itshape,
  numberstyle=\tiny\color{codegray},
  numbers=left,
  numbersep=8pt,
  breaklines=true,
  showstringspaces=false,
  frame=single,
  rulecolor=\color{codegray!40},
  tabsize=2,
  captionpos=b,
  literate={á}{{\'a}}1 {é}{{\'e}}1 {í}{{\'i}}1 {ó}{{\'o}}1 {ú}{{\'u}}1
           {Á}{{\'A}}1 {É}{{\'E}}1 {Í}{{\'I}}1 {Ó}{{\'O}}1 {Ú}{{\'U}}1
           {ñ}{{\~n}}1 {Ñ}{{\~N}}1 {¿}{{?`}}1 {¡}{{!`}}1
}
\lstset{style=winstyle}

% Comando corto para identificadores en línea
\newcommand{\code}[1]{\texttt{#1}}

\newcolumntype{L}[1]{>{\raggedright\arraybackslash}p{#1}}

% ============================================================
%  Portada
% ============================================================
\begin{document}

\begin{titlepage}
  \centering
  \vspace*{2cm}
  {\color{winorange}\rule{\linewidth}{2pt}}\\[0.6cm]
  {\Huge\bfseries WinSports\par}
  \vspace{0.3cm}
  {\LARGE Plataforma de Analítica de Streaming Deportivo\par}
  \vspace{0.2cm}
  {\large Procesamiento de eventos de reproducción sobre MongoDB Atlas,\\
   API en FastAPI y dashboard en React\par}
  {\color{winorange}\rule{\linewidth}{2pt}}\\[1.5cm]

  {\large\bfseries Informe Final\par}
  \vspace{0.4cm}
  {\large Ingeniería de Datos\par}
  \vspace{2.5cm}

  \begin{tabular}{rl}
    \textbf{Autor:} & Juan José Castillo \\[4pt]
    \textbf{Stack:} & MongoDB Atlas $\cdot$ FastAPI $\cdot$ React 19 \\[4pt]
    \textbf{Versión:} & 0.1.0 \\[4pt]
    \textbf{Fecha:} & Mayo de 2026 \\
  \end{tabular}

  \vfill
  {\color{codegray}\small Win Sports — plataforma colombiana de streaming deportivo\par}
\end{titlepage}

\tableofcontents
\newpage

% ============================================================
\section{Resumen}
% ============================================================
WinSports es un sistema de analítica construido sobre los eventos de
reproducción de \textbf{Win Sports}, plataforma colombiana de streaming
deportivo. El sistema ingiere registros crudos de eventos
(\code{START}, \code{BUFFERING}, \code{PAUSE}, \code{COMPLETE}, etc.) desde
archivos CSV hacia \textbf{MongoDB Atlas}, los transforma en colecciones
derivadas mediante un pipeline de agregación, y los expone a través de una
\textbf{API REST en FastAPI} consumida por un \textbf{dashboard interactivo en
React}.

El núcleo del diseño es un modelo de datos de tres niveles
(\code{events} $\rightarrow$ \code{sessions} $\rightarrow$ \code{content\_stats})
que separa la fuente de verdad inmutable de las vistas precalculadas que
alimentan el tablero. Sobre esa base, el dashboard ofrece cuatro perspectivas
analíticas —\emph{Overview}, \emph{QoE} (calidad de experiencia),
\emph{Usuarios} y \emph{Admin}—, además de un módulo de \emph{Alertas}
operativas. El proyecto incluye una suite de pruebas en backend (pytest con
Mongo en memoria) y frontend (Vitest + Testing Library).

% ============================================================
\section{Introducción}
% ============================================================
\subsection{Contexto y motivación}
Cada vez que un usuario reproduce contenido en una plataforma de streaming, el
reproductor emite una secuencia de eventos: el inicio de la reproducción, los
episodios de \emph{buffering}, las pausas, los hitos de progreso (primer
cuartil, mitad, tercer cuartil) y la finalización. Analizar este flujo permite
responder preguntas de negocio clave: ¿qué contenido se consume más?,
¿qué tan buena es la calidad de reproducción?, ¿en qué punto los usuarios
abandonan un contenido?, ¿qué dispositivos predominan?

Estos datos presentan dos características que condicionan el diseño:

\begin{itemize}[leftmargin=*]
  \item \textbf{Esquema flexible.} No todos los eventos traen los mismos campos:
        \code{buffer\_time} solo aparece en eventos de buffering;
        \code{episode} y \code{season} solo en series. Forzar un esquema rígido
        relacional generaría una proliferación de valores nulos.
  \item \textbf{Alto volumen.} En producción, los eventos llegarían como flujo
        en vivo, alcanzando un orden de magnitud de $\sim$100M de documentos.
        La arquitectura debe escalar horizontalmente sin rediseñarse.
\end{itemize}

Por estas razones se eligió \textbf{MongoDB}, una base de datos documental que
maneja naturalmente el esquema flexible y escala horizontalmente vía
\emph{sharding}.

\subsection{Objetivos}
\begin{enumerate}[leftmargin=*]
  \item Diseñar un modelo de datos documental que separe la captura cruda de
        eventos de las métricas analíticas precalculadas.
  \item Implementar un pipeline reproducible que transforme un CSV de eventos
        en colecciones derivadas listas para consulta.
  \item Exponer las métricas mediante una API REST con contratos bien definidos
        y validados.
  \item Construir un dashboard que comunique las métricas de forma visual e
        interactiva a un equipo no técnico.
  \item Garantizar la idempotencia de la carga y la corrección mediante pruebas
        automatizadas.
\end{enumerate}

% ============================================================
\section{Metodología y requisitos}
% ============================================================
\subsection{Enfoque de desarrollo}
El sistema se construyó de forma \textbf{iterativa e incremental}, en cortes
verticales: cada métrica del dashboard se llevó de extremo a extremo (índice y
consulta en la base de datos $\rightarrow$ repositorio $\rightarrow$ esquema y
ruta de la API $\rightarrow$ componente de React) antes de pasar a la siguiente.
Este enfoque mantuvo el sistema funcional en todo momento y permitió validar el
contrato entre capas de forma temprana. El control de versiones se gestionó con
Git, y las dependencias con \code{uv} (backend) y \code{npm} (frontend).

\subsection{Requisitos funcionales}
\begin{center}
\begin{tabular}{@{}lL{10cm}@{}}
\toprule
\textbf{ID} & \textbf{Requisito} \\
\midrule
RF1 & Ingerir un CSV de eventos y persistirlo en la colección \code{events}. \\
RF2 & Derivar automáticamente \code{sessions} y \code{content\_stats} a partir de los eventos. \\
RF3 & Exponer métricas de portada (reproducciones, usuarios, contenido, dispositivos, mapa). \\
RF4 & Exponer métricas de calidad de experiencia (buffer, re-buffering, tiempo de inicio). \\
RF5 & Exponer métricas de comportamiento (funnel, heatmap, perfiles, completitud). \\
RF6 & Detectar y reportar alertas operativas por severidad. \\
RF7 & Permitir la gestión de datos (carga, conteo, edición y borrado) desde un panel. \\
\bottomrule
\end{tabular}
\end{center}

\subsection{Requisitos no funcionales}
\begin{center}
\begin{tabular}{@{}lL{10cm}@{}}
\toprule
\textbf{ID} & \textbf{Requisito} \\
\midrule
RNF1 & \textbf{Idempotencia}: re-cargar el mismo CSV no debe duplicar datos. \\
RNF2 & \textbf{Escalabilidad horizontal} hacia el orden de $\sim$100M de documentos. \\
RNF3 & \textbf{Tolerancia a fallos} en la carga masiva (un documento inválido no detiene el lote). \\
RNF4 & \textbf{Modularidad}: la API nunca accede a Mongo directamente; el front solo consume REST. \\
RNF5 & \textbf{Privacidad}: no exponer identificadores completos de suscriptor en respuestas públicas. \\
RNF6 & \textbf{Respuesta interactiva}: consultas del dashboard servidas desde vistas precalculadas. \\
RNF7 & \textbf{Verificabilidad}: cobertura de pruebas en backend y frontend. \\
\bottomrule
\end{tabular}
\end{center}

% ============================================================
\section{Arquitectura del sistema}
% ============================================================
El sistema se organiza en capas con responsabilidades estrictamente separadas.
La API nunca accede a MongoDB directamente: toda consulta pasa por la capa de
repositorios. El frontend solo consume la API REST. Esta modularidad permite que
base de datos, API y frontend evolucionen de forma independiente.

\subsection{Stack tecnológico}

\begin{center}
\begin{tabular}{@{}ll@{}}
\toprule
\textbf{Capa} & \textbf{Tecnología} \\
\midrule
Base de datos & MongoDB Atlas + Motor (driver asíncrono) \\
API & FastAPI + Uvicorn + Pydantic v2 \\
Frontend & React 19 + Vite + Recharts + Leaflet \\
Entorno backend & Python 3.14 + uv \\
Entorno frontend & Node.js 18+ \\
Pruebas & pytest + mongomock-motor $\cdot$ Vitest + Testing Library \\
\bottomrule
\end{tabular}
\end{center}

\subsection{Justificación: NoSQL frente a modelo relacional}
La elección de una base de datos documental no es arbitraria; responde a la
naturaleza del dato y a los requisitos no funcionales:

\begin{center}
\begin{tabular}{@{}L{3.2cm}L{5.3cm}L{5.3cm}@{}}
\toprule
\textbf{Criterio} & \textbf{Relacional (SQL)} & \textbf{Documental (MongoDB)} \\
\midrule
Esquema & Rígido; campos opcionales obligan a tablas dispersas o muchos \code{NULL}. & Flexible; cada evento lleva solo los campos que aplican. \\
Escalado & Vertical, principalmente; \emph{sharding} complejo. & Horizontal nativo vía \emph{sharding}. \\
Lectura analítica & \emph{Joins} costosos sobre tablas grandes. & Vistas materializadas (colecciones derivadas) sin \emph{joins}. \\
Ingesta masiva & Validación estricta detiene lotes ante un error. & \code{validationAction: warn} permite degradar con gracia. \\
Sincronización & Requiere triggers / ETL externo. & \emph{Change streams} nativos. \\
\bottomrule
\end{tabular}
\end{center}

\subsection{Separación de capas}
La estructura del repositorio refleja directamente la arquitectura:

\begin{lstlisting}[language=bash,caption={Organización por responsabilidades}]
db/collections/   validators + indices + setup de cada coleccion
db/indexes/       definicion de indices (importados por collections)
db/pipelines/     logica de construccion de colecciones derivadas
db/repositories/  queries hacia Mongo (una funcion por operacion)
db/csv_ingest.py  parseo/ingesta de CSV (compartido por script y admin)
api/schemas/      modelos Pydantic - forma del JSON que recibe React
api/routes/       endpoints FastAPI - une repositories con schemas
config/           settings (.env) + constantes (colecciones, centroides)
frontend/         React + Vite (pages, components, layout, styles, test)
\end{lstlisting}

\subsection{Flujo de datos general}
El dato recorre el sistema en un solo sentido, de la fuente cruda hacia las
vistas agregadas:

\begin{center}
\fbox{\code{CSV}} $\rightarrow$ \fbox{\code{load\_csv.py}} $\rightarrow$
\fbox{\code{events}} $\xrightarrow{\text{change stream}}$
\fbox{\code{sessions}} $\xrightarrow{\text{change stream}}$
\fbox{\code{content\_stats}} $\rightarrow$ \fbox{API} $\rightarrow$ \fbox{React}
\end{center}

En \textbf{desarrollo}, \code{load\_csv.py} inserta los eventos y
\code{scripts/test\_pipeline.py} construye manualmente las colecciones
derivadas. En \textbf{producción}, los \emph{change streams} de MongoDB mantienen
\code{sessions} y \code{content\_stats} sincronizadas en tiempo real ante cada
nuevo evento.

% ============================================================
\section{Conjunto de datos}
% ============================================================
\subsection{Origen y volumen}
La fuente es un CSV de eventos de reproducción de Win Sports. El repositorio
incluye dos archivos de muestra para desarrollo y pruebas:
\code{data/50\_data.csv} (50 eventos) y \code{data/1000\_data.csv} (1.000
eventos). El pipeline se validó además con cargas del orden de
\textbf{cientos de miles de eventos} ($\sim$300k), confirmando que el modelo de
colecciones derivadas sostiene el dashboard sin degradar los tiempos de
respuesta.

\subsection{Diccionario de datos}
El CSV original trae 22 columnas. La función \code{parse\_row} en
\code{db/csv\_ingest.py} las mapea a los campos del documento de \code{events},
casteando cada una al \code{bsonType} esperado:

\begin{center}
\begin{longtable}{@{}lll@{}}
\toprule
\textbf{Columna CSV} & \textbf{Campo evento} & \textbf{Tipo} \\
\midrule
\endhead
\code{Date} & \code{date} & date \\
\code{CountryCode} & \code{country\_code} & string \\
\code{SubscriberID} & \code{subscriber\_id} & string \\
\code{CustomerId} & \code{customer\_id} & string \\
\code{ContentType} & \code{content\_type} & string \\
\code{Title} & \code{title} & string \\
\code{Episode} & \code{episode} & int \\
\code{SeriesTitle} & \code{series\_title} & string \\
\code{ReleaseYear} & \code{release\_year} & int \\
\code{Duration} & \code{duration} & int \\
\code{Season} & \code{season} & int \\
\code{Genres} & \code{genres} & string \\
\code{DeviceType} & \code{device\_type} & string (enum) \\
\code{TypeEvent} & \code{type\_event} & string \\
\code{Position} & \code{position} & int \\
\code{Language} & \code{language} & string \\
\code{Bitrate} & \code{bitrate} & int \\
\code{BufferTime} & \code{buffer\_time} & double \\
\code{PlaybackNetTime} & \code{playback\_net\_time} & int \\
\code{deviceDescription} & \code{device\_description} & string \\
\code{CALC\_ProgramType} & \code{calc\_program\_type} & string \\
\code{CALC\_BitrateType} & \code{calc\_bitrate\_type} & string \\
\bottomrule
\end{longtable}
\end{center}

\subsection{Catálogo de tipos de evento}
El campo \code{type\_event} dirige toda la lógica de negocio. Los tipos
relevantes son:

\begin{center}
\begin{tabular}{@{}lL{8.5cm}@{}}
\toprule
\textbf{Tipo} & \textbf{Significado} \\
\midrule
\code{START} & Inicia una reproducción (delimita el comienzo de una sesión). \\
\code{START-BUFFERING} & Buffer previo al primer frame (mide el \emph{startup time}). \\
\code{RE-BUFFERING} & Interrupción por buffer durante la reproducción. \\
\code{PAUSE} & Pausa del usuario. \\
\code{FIRSTQUARTILE} & Se alcanzó el 25\% del contenido. \\
\code{MIDPOINT} & Se alcanzó el 50\% del contenido. \\
\code{THIRDQUARTILE} & Se alcanzó el 75\% del contenido. \\
\code{COMPLETE} & El contenido se vio hasta el final. \\
\bottomrule
\end{tabular}
\end{center}

\subsection{Ejemplo de filas crudas}
\begin{lstlisting}[language=bash,caption={Dos eventos del CSV (cabecera abreviada)}]
Date,CountryCode,SubscriberID,CustomerId,...,DeviceType,TypeEvent,Position,...,BufferTime,...
2026-01-01 08:58:00,CO,0d16...20a9,6955...de75,...,WEB,START,0,...,,...
2026-01-01 08:58:05,CO,0d16...20a9,6955...de75,...,WEB,START-BUFFERING,0,...,6.805,...
\end{lstlisting}

Nótese que \code{BufferTime} solo viene poblado en el evento de buffering: ese es
exactamente el tipo de esquema flexible que motiva el uso de MongoDB.

% ============================================================
\section{Modelo de datos}
% ============================================================
El modelo se compone de tres colecciones con responsabilidades separadas. Cada
colección downstream es una vista materializada de la anterior, lo que evita
recalcular agregaciones costosas en cada petición del dashboard.

\subsection{Colección \code{events} — fuente de verdad}
Almacena un documento por cada evento del CSV. Es un \emph{log} crudo,
inmutable: nunca se modifica tras la inserción. Su validador
(\code{\$jsonSchema}) exige cinco campos obligatorios y deja el resto
opcional, respetando el esquema flexible de los eventos de streaming.

\begin{center}
\begin{tabular}{@{}lL{8.5cm}@{}}
\toprule
\textbf{Campo} & \textbf{Descripción} \\
\midrule
\code{date} & Fecha-hora del evento (obligatorio, \code{bsonType: date}). \\
\code{subscriber\_id} & Identificador del suscriptor (obligatorio). \\
\code{customer\_id} & Identificador de la pieza de contenido (obligatorio). \\
\code{type\_event} & Tipo de evento: \code{START}, \code{PAUSE}, \code{COMPLETE}, etc. (obligatorio). \\
\code{device\_type} & Plataforma: \code{ANDR}, \code{IOS} o \code{WEB} (obligatorio, \emph{enum}). \\
\code{buffer\_time} & Tiempo de buffer; solo en eventos de buffering. \\
\code{episode}, \code{season} & Solo en contenido tipo episodio/serie. \\
\code{position} & Posición de reproducción, en segundos. \\
\code{country\_code} & Código de país (alimenta el mapa de usuarios). \\
\bottomrule
\end{tabular}
\end{center}

\paragraph{Validación.} El validador usa \code{validationAction: "warn"} en
lugar de \code{"error"}: durante una carga masiva, un documento con un campo
inesperado no debe detener la inserción de millones de filas; el problema se
registra sin interrumpir el proceso. Los enteros aceptan
\code{["int", "long"]} porque \code{bitrate} y \code{position} pueden superar el
rango de \code{int32} en el dataset completo.

\begin{lstlisting}[language=Python,caption={Documento de ejemplo en \code{events}}]
{
  "date": ISODate("2026-01-01T08:58:05Z"),
  "country_code": "CO",
  "subscriber_id": "0d1687a8...390020a9",
  "customer_id": "69553fe8c788eafebe63de75",
  "content_type": "MOVIE",
  "title": "PARTIDO APARTE",
  "duration": 3540,
  "device_type": "WEB",
  "type_event": "START-BUFFERING",
  "position": 0,
  "buffer_time": 6.805,
  "calc_bitrate_type": "HIGH"
}
\end{lstlisting}

\subsection{Colección \code{sessions} — derivada de events}
Una \emph{sesión} agrupa todos los eventos de un par
(\code{subscriber\_id}, \code{customer\_id}) entre un \code{START} y el
siguiente \code{START}. Precalcula las métricas de progreso y calidad para no
recalcularlas sobre el conjunto completo de eventos en cada consulta.

\begin{itemize}[leftmargin=*]
  \item \textbf{Clave única:} \code{subscriber\_id + customer\_id + start\_time}.
  \item \textbf{Progreso:} \code{reached\_firstquartile / midpoint /
        thirdquartile / complete} (booleanos del funnel), \code{completion\_pct},
        \code{max\_position}.
  \item \textbf{Calidad (QoE):} \code{total\_buffer\_time}, \code{rebuffer\_count},
        \code{pause\_count}, \code{dominant\_bitrate\_type}.
\end{itemize}

\begin{lstlisting}[language=Python,caption={Documento de ejemplo en \code{sessions}}]
{
  "subscriber_id": "0d1687a8...390020a9",
  "customer_id": "69553fe8c788eafebe63de75",
  "start_time": ISODate("2026-01-01T08:58:00Z"),
  "end_time": ISODate("2026-01-01T09:55:00Z"),
  "device_type": "WEB",
  "duration": 3540,
  "total_buffer_time": 6.805,
  "rebuffer_count": 0,
  "pause_count": 1,
  "max_position": 3120,
  "completion_pct": 88.14,
  "reached_firstquartile": true,
  "reached_midpoint": true,
  "reached_thirdquartile": true,
  "reached_complete": false,
  "event_count": 7
}
\end{lstlisting}

\subsection{Colección \code{content\_stats} — derivada de sessions}
Un documento por \code{customer\_id} (pieza de contenido), con métricas
agregadas que alimentan directamente los rankings del dashboard.
Importante: \code{customer\_id} representa una pieza concreta (un episodio, una
película, un partido), no un programa completo; para métricas por serie se
agrupa por \code{series\_title}.

\begin{itemize}[leftmargin=*]
  \item \textbf{Volumen:} \code{total\_plays}, \code{unique\_viewers}.
  \item \textbf{Funnel:} \code{firstquartile\_rate}, \code{midpoint\_rate},
        \code{thirdquartile\_rate}, \code{completion\_rate}.
  \item \textbf{Calidad:} \code{avg\_buffer\_time}, \code{rebuffer\_rate},
        \code{avg\_completion\_pct}.
  \item \textbf{Breakdown:} \code{plays\_by\_device} $=
        \{\code{ANDR}, \code{IOS}, \code{WEB}\}$.
\end{itemize}

\begin{lstlisting}[language=Python,caption={Documento de ejemplo en \code{content\_stats}}]
{
  "customer_id": "69553fe8c788eafebe63de75",
  "title": "PARTIDO APARTE",
  "content_type": "MOVIE",
  "total_plays": 40,
  "unique_viewers": 33,
  "firstquartile_rate": 82.5,
  "midpoint_rate": 67.5,
  "thirdquartile_rate": 50.0,
  "completion_rate": 30.0,
  "avg_buffer_time": 2.314,
  "rebuffer_rate": 18.5,
  "avg_completion_pct": 64.2,
  "plays_by_device": { "WEB": 22, "ANDR": 12, "IOS": 6 }
}
\end{lstlisting}

\subsection{Estrategia de índices}
Cada colección define índices alineados con sus patrones de consulta. Por
ejemplo, en \code{events} el índice compuesto
\code{(subscriber\_id, customer\_id, date)} acelera la reconstrucción de
sesiones, mientras que \code{(type\_event, date DESC)} sirve a los filtros de
QoE. En \code{sessions}, el índice único
\code{(subscriber\_id, customer\_id, start\_time)} es la clave de upsert que
garantiza idempotencia.

\paragraph{Índices de \code{events}.}
\begin{center}
\begin{tabular}{@{}lL{7.5cm}@{}}
\toprule
\textbf{Índice} & \textbf{Propósito} \\
\midrule
\code{(subscriber\_id, customer\_id, date)} & Reconstrucción ordenada de sesiones. \\
\code{(date DESC)} & Ventanas de tiempo (alertas, actividad reciente). \\
\code{(type\_event, date DESC)} & Filtros de QoE por tipo de evento. \\
\code{(customer\_id, type\_event)} & Rankings de contenido. \\
\code{(device\_type)} & Desglose por plataforma. \\
\bottomrule
\end{tabular}
\end{center}

\paragraph{Índices de \code{sessions}.}
\begin{center}
\begin{tabular}{@{}lL{7.5cm}@{}}
\toprule
\textbf{Índice} & \textbf{Propósito} \\
\midrule
\code{(subscriber\_id)} & Usuarios únicos y perfil de usuario. \\
\code{(customer\_id)} & Consultas por contenido. \\
\code{(device\_type)} & Ranking de dispositivos. \\
\code{(start\_time DESC, total\_buffer\_time DESC)} & Alertas de buffer reciente. \\
\code{(customer\_id, reached\_complete)} & Funnel de retención por contenido. \\
\code{(subscriber\_id, customer\_id, start\_time)} \textbf{UNIQUE} & Clave de upsert (idempotencia). \\
\bottomrule
\end{tabular}
\end{center}

\paragraph{Índices de \code{content\_stats}.}
\begin{center}
\begin{tabular}{@{}lL{7.5cm}@{}}
\toprule
\textbf{Índice} & \textbf{Propósito} \\
\midrule
\code{(customer\_id)} \textbf{UNIQUE} & Un documento por contenido. \\
\code{(total\_plays DESC)} & Ranking por número de vistas. \\
\code{(avg\_buffer\_time DESC)} & Ranking QoE (peores por buffer). \\
\code{(genres, total\_plays DESC)} & Rankings filtrados por género. \\
\code{(series\_title, total\_plays DESC)} & Agrupación por serie. \\
\bottomrule
\end{tabular}
\end{center}

% ============================================================
\section{Ingesta y pipeline de datos}
% ============================================================
\subsection{Parseo e ingesta de CSV}
La lógica de parseo vive en \code{db/csv\_ingest.py} y es compartida por el
script de línea de comandos (\code{scripts/load\_csv.py}) y el endpoint del panel
de administración (\code{POST /api/admin/upload-csv}). Así, el script y la UI
procesan los CSV de forma idéntica. El módulo realiza:

\begin{enumerate}[leftmargin=*]
  \item Casteo de tipos: cada columna del CSV se mapea al \code{bsonType}
        esperado. Los \code{NaN} de pandas se convierten en \code{None} y los
        enteros que vienen como \code{"2023.0"} se castean vía \code{float}.
  \item Filtrado: se descarta toda fila sin alguno de los cinco campos
        obligatorios (no se puede construir una sesión a partir de ella).
  \item Inserción por lotes de 1.000 documentos con \code{ordered=False}.
\end{enumerate}

\paragraph{Tolerancia a fallos.} Con \code{ordered=False}, si un documento falla
la validación de Mongo, el lote continúa en lugar de detenerse. El conteo real
de insertados se toma del resultado de cada lote, incluyendo el \code{nInserted}
que reporta un \code{BulkWriteError}.

\subsection{Construcción de sesiones}
El pipeline de \code{db/pipelines/sessions.py} agrupa los eventos por par
(\code{subscriber\_id}, \code{customer\_id}), los ordena por fecha y los parte en
sub-sesiones cada vez que aparece un \code{START}. Por cada sub-sesión calcula
el buffer total, los conteos de re-buffering y pausa, la posición máxima y el
porcentaje de completitud.

\begin{lstlisting}[language=Python,caption={Partición en sub-sesiones por evento START (db/pipelines/sessions.py)}]
for ev in evs:
    if ev["type_event"] == "START" and current:
        sub_sessions.append(current)
        current = []
    current.append(ev)
if current:
    sub_sessions.append(current)
\end{lstlisting}

\paragraph{Cap de \code{completion\_pct} a 100\%.} Algunas filas traen
\code{Position} en una unidad distinta a \code{Duration} (milisegundos frente a
segundos) o posiciones residuales de sesiones previas, lo que dispara el
porcentaje por encima de 100. Capear a 100\% evita que esos valores atípicos
contaminen los promedios de completitud del dashboard:

\begin{lstlisting}[language=Python,caption={Cálculo de completitud capado}]
completion_pct = (
    round(min(max_position / duration * 100, 100.0), 2)
    if duration and duration > 0 and max_position is not None
    else None
)
\end{lstlisting}

\subsection{Agregación de content\_stats}
El pipeline de \code{db/pipelines/content\_stats.py} lee todas las sesiones, las
agrupa por \code{customer\_id} y produce un documento agregado por contenido:
volumen, tasas del funnel, métricas de calidad y desglose por dispositivo. La
escritura usa \code{UpdateOne(..., upsert=True)} por \code{customer\_id}.

\subsection{Idempotencia mediante upsert}
Tanto sesiones como content\_stats se escriben con \emph{upsert} sobre su clave
única. Cargar el mismo CSV dos veces \textbf{actualiza} los documentos en lugar
de duplicarlos, lo que hace que el pipeline sea seguro de re-ejecutar.

% ============================================================
\section{API REST}
% ============================================================
Toda la API vive bajo el prefijo \code{/api}. Las respuestas son JSON; los
errores retornan \code{\{ "detail": "mensaje en español" \}}. Los contratos se
validan con esquemas Pydantic v2 y la documentación interactiva (Swagger) se
genera automáticamente en \code{/docs}.

\subsection{Módulo Overview}
\begin{center}
\begin{tabular}{@{}lL{7.5cm}@{}}
\toprule
\textbf{Endpoint} & \textbf{Descripción} \\
\midrule
\code{GET /overview/unique-users} & Total de suscriptores distintos. \\
\code{GET /overview/total-plays} & Total de reproducciones (= sesiones). \\
\code{GET /overview/device-ranking} & Sesiones agrupadas por dispositivo. \\
\code{GET /overview/activity-by-hour} & Sesiones por hora del día (las 24 horas). \\
\code{GET /overview/top-content} & Contenidos más vistos (param \code{limit}). \\
\code{GET /overview/users-map} & Una burbuja por usuario activo, ubicada en el centroide de su país + jitter determinístico. \\
\bottomrule
\end{tabular}
\end{center}

El endpoint del mapa trunca el \code{subscriber\_id} a 8 caracteres (no expone
el identificador completo) y posiciona cada usuario en el centroide de su país
más un \emph{jitter} determinístico, ya que el dataset solo trae el código de
país, no coordenadas.

\subsection{Módulo QoE}
\begin{center}
\begin{tabular}{@{}lL{7.5cm}@{}}
\toprule
\textbf{Endpoint} & \textbf{Descripción} \\
\midrule
\code{GET /qoe/buffer-by-content} & Buffer promedio por contenido (peores primero). \\
\code{GET /qoe/rebuffering-rate} & Tasa de re-buffering por eventos y por sesiones. \\
\code{GET /qoe/startup-time} & Tiempo de inicialización (promedio, máximo y distribución). \\
\code{GET /qoe/event-ranking} & Ranking de tipos de evento por frecuencia. \\
\bottomrule
\end{tabular}
\end{center}

\subsection{Módulo Users}
\begin{center}
\begin{tabular}{@{}lL{7.5cm}@{}}
\toprule
\textbf{Endpoint} & \textbf{Descripción} \\
\midrule
\code{GET /users/retention-funnel} & Funnel: cuántas sesiones llegaron a cada cuartil. \\
\code{GET /users/activity-heatmap} & Sesiones por hora $\times$ día de la semana (168 puntos). \\
\code{GET /users/content-completion-ranking} & Contenidos más completados y más abandonados (param \code{min\_plays}). \\
\code{GET /users/user-profiles} & Clasificación de usuarios por actividad. \\
\bottomrule
\end{tabular}
\end{center}

\paragraph{Perfiles de usuario.} Los usuarios se clasifican por el número total
de eventos que generan: \textbf{Vague} ($<10$), \textbf{Mid} ($10$–$49$) y
\textbf{Heavy} ($\geq 50$). Los umbrales son constantes
(\code{PROFILE\_MID\_MIN}, \code{PROFILE\_HEAVY\_MIN}) en
\code{db/repositories/users.py}, y la clasificación se resuelve del lado del
servidor con un \code{\$switch} en la agregación.

\subsection{Módulo Admin}
\begin{center}
\begin{tabular}{@{}lL{7.5cm}@{}}
\toprule
\textbf{Endpoint} & \textbf{Descripción} \\
\midrule

\code{GET /admin/collections} & Conteo de documentos por colección. \\
\code{POST /admin/upload-csv} & Sube un CSV y re-construye las derivadas. \\
\code{GET /admin/documents/\{col\}} & Documentos paginados (\code{page}, \code{page\_size}). \\
\code{PUT /admin/documents/\{col\}/\{id\}} & Edición en vivo (\code{\$set} de campos). \\
\code{DELETE /admin/documents/\{col\}/\{id\}} & Borrado de un documento. \\
\bottomrule
\end{tabular}
\end{center}

\subsection{Ejemplos de petición y respuesta}
Por ejemplo, el ranking de contenido devuelve un arreglo ya ordenado y
proyectado:

\begin{lstlisting}[language=bash,caption={Petición y respuesta de top-content}]
GET /api/overview/top-content?limit=2

{
  "content": [
    { "title": "Liga BetPlay", "total_plays": 120,
      "unique_viewers": 80, "content_type": "LIVE" },
    { "title": "PARTIDO APARTE", "total_plays": 40,
      "unique_viewers": 33, "content_type": "MOVIE" }
  ]
}
\end{lstlisting}

\begin{lstlisting}[language=bash,caption={Respuesta del funnel de retención}]
GET /api/users/retention-funnel

{ "total": 1000, "firstquartile": 800, "midpoint": 600,
  "thirdquartile": 400, "complete": 200 }
\end{lstlisting}

\paragraph{Las consultas se resuelven en el servidor.} Ningún repositorio trae
documentos completos: las métricas se calculan con el \emph{aggregation
framework} de MongoDB. La clasificación de perfiles, por ejemplo, usa
\code{\$switch} para etiquetar a cada usuario sin traer datos al cliente:

\begin{lstlisting}[language=Python,caption={Clasificación de perfiles con \$switch (db/repositories/users.py)}]
"profile": {
  "$switch": {
    "branches": [
      {"case": {"$lt": ["$events", PROFILE_MID_MIN]},   "then": "Vague"},
      {"case": {"$lt": ["$events", PROFILE_HEAVY_MIN]},  "then": "Mid"},
    ],
    "default": "Heavy",
  }
}
\end{lstlisting}

\subsection{Manejo de errores}
\begin{center}
\begin{tabular}{@{}cl@{}}
\toprule
\textbf{Código} & \textbf{Significado} \\
\midrule
\code{400} & Petición inválida (ID malformado, archivo no-CSV o vacío). \\
\code{404} & Recurso o colección no encontrada. \\
\code{422} & Error de validación de parámetros / cuerpo. \\
\code{500} & Error interno del servidor. \\
\bottomrule
\end{tabular}
\end{center}

% ============================================================
\section{Dashboard (Frontend)}
% ============================================================
El frontend es una SPA en \textbf{React 19} servida con \textbf{Vite}. Las
gráficas se construyen con \textbf{Recharts} y el mapa de usuarios con
\textbf{Leaflet} (\code{react-leaflet}). El proxy de Vite redirige \code{/api}
hacia \code{http://localhost:8000}, por lo que no se requiere configurar CORS en
desarrollo.

\subsection{Páginas}
\begin{center}
\begin{tabular}{@{}lL{7cm}c@{}}
\toprule
\textbf{Página} & \textbf{Contenido} & \textbf{Estado} \\
\midrule
Overview & KPIs, mapa de burbujas, ranking de dispositivos, actividad por hora & Listo \\
QoE & Buffer por contenido, tasa de re-buffering, tiempo de inicio, ranking de eventos & Listo \\
Usuarios & Perfiles, funnel de retención, heatmap, completitud por contenido & Listo \\
Admin & Carga de CSV, conteo de colecciones, edición/borrado en vivo & Listo \\
Alertas & Tablero de alertas operativas & En construcción \\
\bottomrule
\end{tabular}
\end{center}

\subsection{Componentes y temas}
Cada página orquesta componentes presentacionales que consumen un endpoint
específico (por ejemplo, \code{KPICard}, \code{UsersMap}, \code{RebufferingRate},
\code{RetentionFunnel}, \code{DocumentTable}). El dashboard define dos temas en
\code{src/styles/theme.css}: \textbf{brand} (naranja Win Sports \code{\#FF6B00}
sobre gris oscuro) y \textbf{premium} (naranja quemado sobre casi negro).

\subsection{Mapa de usuarios: jitter determinístico}
Como el dataset solo trae el código de país (no coordenadas), cada usuario se
ubica en el centroide de su país más un desplazamiento (\emph{jitter}) calculado
de forma determinística a partir del hash MD5 de su \code{subscriber\_id}. Así,
un mismo usuario cae siempre en el mismo punto entre peticiones, evitando que
todos se apilen exactamente sobre el centroide:

\begin{lstlisting}[language=Python,caption={Jitter determinístico (db/repositories/overview.py)}]
h = hashlib.md5(seed.encode()).hexdigest()
lat_off = (int(h[:8], 16) / 0xFFFFFFFF - 0.5) * 2 * spread
lon_off = (int(h[8:16], 16) / 0xFFFFFFFF - 0.5) * 2 * spread
\end{lstlisting}

% ============================================================
\section{Sistema de alertas}
% ============================================================
El módulo de alertas (\code{db/repositories/alerts.py}) detecta condiciones
operativas anómalas dentro de la ventana temporal de los datos cargados,
clasificándolas en tres niveles de severidad:

\begin{center}
\begin{tabular}{@{}clL{6.5cm}@{}}
\toprule
\textbf{Severidad} & \textbf{Tipo} & \textbf{Condición} \\
\midrule
Roja & \code{buffer\_critico} & Sesiones con más de 15 s de buffer acumulado. \\
Amarilla & \code{rebuffer\_rate} & Contenidos donde $>30\%$ de las sesiones (mín. 10) tienen al menos un re-buffering. \\
Azul & \code{pause\_storm} & Usuarios con más de 5 pausas en menos de 8 minutos en una misma sesión. \\
\bottomrule
\end{tabular}
\end{center}

El endpoint \code{GET /api/alerts/} combina las tres familias, las ordena por
severidad y marca temporal, y retorna además los totales por color y la ventana
del período analizado. La lógica del backend está implementada; la página de
Alertas del dashboard se encuentra en construcción.

% ============================================================
\section{Panel de administración}
% ============================================================
La pestaña \textbf{Admin} permite al equipo gestionar los datos sin tocar Mongo
directamente:

\begin{itemize}[leftmargin=*]
  \item \textbf{Carga de CSV.} Sube un archivo, lo inserta en \code{events}
        reutilizando \code{db/csv\_ingest.py}, y re-construye \code{sessions} y
        \code{content\_stats}. Muestra un resumen: insertados, omitidos, total y
        cuántas derivadas quedaron.
  \item \textbf{Conteo de colecciones.} Tarjetas con el total y el desglose por
        colección, refrescadas tras cada operación.
  \item \textbf{Edición/borrado en vivo.} Tabla paginada donde se editan campos
        escalares (preservando su tipo original) o se borran documentos.
\end{itemize}

\paragraph{Seguridad de la capa de admin.} Solo se exponen tres colecciones
(\code{events}, \code{sessions}, \code{content\_stats}); cualquier otra se
rechaza con \code{404}, de modo que el panel no puede tocar colecciones internas
de Mongo. El \code{\_id} es inmutable y los \code{ObjectId} se validan antes de
operar (un ID malformado retorna \code{400}).

\paragraph{Advertencia.} El panel \textbf{aún no tiene autenticación}; asume que
la API corre en una red interna. Añadir auth es un pendiente antes de exponerlo a
producción. Editar una colección derivada es un cambio manual que una nueva carga
de CSV sobrescribe; para cambios permanentes hay que editar \code{events} y
re-procesar.

% ============================================================
\section{Decisiones de diseño}
% ============================================================
\begin{description}[leftmargin=0pt,style=nextline]
  \item[Change streams en vez de recalcular por request.] Con millones de
        eventos, agregar en tiempo real sería demasiado lento para un dashboard
        interactivo. Las colecciones derivadas actúan como una caché persistente
        que los change streams mantienen al día.
  \item[\code{validationAction: "warn"} en vez de \code{"error"}.] Durante una
        carga masiva, un documento con un campo inesperado no debe detener la
        inserción del resto. Los errores se registran sin interrumpir.
  \item[Upsert con clave única en sessions.]
        \code{subscriber\_id + customer\_id + start\_time} garantiza
        idempotencia: re-cargar el mismo CSV no duplica sesiones.
  \item[Scripts separados para carga y pipeline.] \code{load\_csv.py} solo
        inserta eventos; \code{test\_pipeline.py} construye las derivadas. Esto
        permite iterar el pipeline sin volver a subir el CSV completo.
  \item[Parseo de CSV compartido.] El script CLI y el panel de admin usan el
        mismo módulo (\code{csv\_ingest.py}), evitando divergencias entre ambos
        caminos de ingesta.
  \item[Degradación con gracia de los change streams.] Si el clúster no soporta
        change streams (Mongo \emph{standalone}, sin réplica), la API loguea el
        problema y sigue funcionando con las derivadas que construye
        \code{test\_pipeline.py}.
\end{description}

% ============================================================
\section{Escalabilidad y requisitos no funcionales}
% ============================================================
El sistema está diseñado para escalar del dataset de prueba al orden de
magnitud de la operación real:

\begin{itemize}[leftmargin=*]
  \item \textbf{Sharding horizontal.} MongoDB particiona \code{events} por su
        clave natural de acceso, repartiendo carga y almacenamiento entre nodos
        sin cambiar la aplicación.
  \item \textbf{Vistas materializadas.} Las colecciones derivadas evitan agregar
        sobre millones de eventos en cada petición; el costo de cómputo se paga
        una vez en la ingesta, no en cada consulta del dashboard.
  \item \textbf{Agregación del lado del servidor.} Todos los repositorios usan
        \code{\$group}/\code{\$project}/\code{count\_documents}; no se traen
        documentos completos al proceso de Python, acotando el uso de memoria.
  \item \textbf{Ingesta por lotes tolerante a fallos.} La inserción en lotes de
        1.000 con \code{ordered=False} sostiene cargas grandes sin abortar ante
        documentos puntuales inválidos.
  \item \textbf{Privacidad.} Las respuestas públicas truncan el identificador de
        suscriptor; el mapa nunca expone el \code{subscriber\_id} completo.
\end{itemize}

% ============================================================
\section{Pruebas}
% ============================================================
\subsection{Backend}
La suite de pytest usa \textbf{mongomock-motor} para simular MongoDB en memoria,
de modo que las pruebas no dependen de un clúster real. Cubre:

\begin{itemize}[leftmargin=*]
  \item La lógica de los pipelines (construcción de sesiones y content\_stats).
  \item El parseo y filtrado del CSV.
  \item Los repositorios de consulta.
  \item El contrato de todas las rutas de la API, incluido el panel de Admin.
\end{itemize}

\begin{lstlisting}[language=bash,caption={Ejecución de pruebas}]
# Backend
uv run pytest

# Frontend
cd frontend && npm run test
\end{lstlisting}

\subsection{Frontend}
Las pruebas usan \textbf{Vitest} + \textbf{Testing Library} sobre \textbf{jsdom}.
El setup global mockea \code{react-leaflet} (jsdom no monta el mapa) y
\code{fetch}. Las pruebas se enfocan en los componentes presentacionales
(estados de carga, vacío y con datos) y en la navegación del \emph{shell}.

% ============================================================
\section{Riesgos y limitaciones}
% ============================================================
\subsection{Resueltos}
\begin{itemize}[leftmargin=*]
  \item \textbf{Duplicados en sessions.} Blindado con el índice único
        \code{idx\_sessions\_unique\_key} y escrituras vía \code{upsert} sobre la
        clave compuesta. Un re-procesamiento actualiza la sesión en vez de
        duplicarla.
  \item \textbf{Consumo de memoria en consultas.} Ningún repositorio trae
        documentos completos: las consultas resuelven con agregaciones
        \code{\$group}/\code{\$project} o \code{count\_documents}, que agregan del
        lado del servidor.
\end{itemize}

\subsection{Abiertos / a vigilar}
\begin{itemize}[leftmargin=*]
  \item \code{distinct("subscriber\_id")} (usado al construir sesiones) tiene un
        tope de 16 MB de BSON. Si el número de suscriptores crece mucho, conviene
        migrar a \code{\$group + \$count}.
  \item El bundle del frontend supera los 500 kB (Leaflet + Recharts). Hoy es
        solo un \emph{warning} de build; si molesta, se puede aplicar
        \emph{code-splitting} con \code{import()} dinámico del mapa.
  \item El panel de Admin carece de autenticación.
\end{itemize}

% ============================================================
\section{Trabajo futuro}
% ============================================================
\begin{itemize}[leftmargin=*]
  \item Completar la página de Alertas del dashboard (el backend ya existe).
  \item Añadir autenticación/autorización al panel de administración.
  \item Editor tipado de fechas y de campos anidados en la tabla de Admin.
  \item Búsqueda/filtro dentro de la tabla de documentos.
  \item Despliegue de los change streams en un clúster réplica para
        sincronización en tiempo real.
  \item Migrar consultas sensibles a volumen (\code{distinct}) a agregaciones.
\end{itemize}

% ============================================================
\section{Conclusiones}
% ============================================================
WinSports demuestra cómo un modelo de datos documental de tres niveles
—fuente de verdad inmutable más vistas materializadas— resuelve de forma
elegante la tensión entre la captura de eventos crudos y la necesidad de
consultas analíticas rápidas. La separación estricta de capas (base de datos,
repositorios, API y frontend) hace que el sistema sea mantenible y extensible:
cada componente puede evolucionar de forma independiente mientras respete el
contrato que lo conecta con el siguiente.

Las decisiones clave —change streams como caché persistente, upsert idempotente,
parseo de CSV compartido, validación tolerante a fallos— responden directamente
a las dos características del dominio: esquema flexible y alto volumen. El
resultado es un sistema que, además de funcionar sobre el dataset de prueba,
está diseñado para escalar al orden de magnitud que exigiría la operación real
de Win Sports.

% ============================================================
\section{Glosario}
% ============================================================
\begin{description}[leftmargin=0pt,style=nextline]
  \item[Evento] Registro atómico emitido por el reproductor (\code{START},
        \code{PAUSE}, \code{COMPLETE}, etc.). Una fila del CSV.
  \item[Sesión] Conjunto de eventos de un par usuario--contenido entre un
        \code{START} y el siguiente. Unidad de análisis de reproducción.
  \item[QoE (Quality of Experience)] Calidad percibida de la reproducción:
        buffer, re-buffering y tiempo de inicialización.
  \item[Funnel de retención] Proporción de sesiones que alcanzan cada hito de
        progreso (primer cuartil, mitad, tercer cuartil, completo).
  \item[Upsert] Operación que actualiza el documento si existe o lo inserta si
        no; base de la idempotencia del pipeline.
  \item[Change stream] Mecanismo de MongoDB que emite cambios en una colección,
        usado para sincronizar las colecciones derivadas en tiempo real.
  \item[Vista materializada] Colección derivada precalculada (\code{sessions},
        \code{content\_stats}) que actúa como caché de consultas.
  \item[Centroide] Coordenada (lat, lon) representativa de un país, usada como
        ancla del mapa de usuarios.
  \item[Jitter] Desplazamiento determinístico aplicado a un punto del mapa para
        que usuarios del mismo país no se solapen.
\end{description}

% ============================================================
\section{Referencias}
% ============================================================
\renewcommand{\refname}{}
\begin{thebibliography}{99}

\bibitem{mongodb} MongoDB, Inc.
  \emph{MongoDB Manual — Documents, Schema Validation y Change Streams}.
  \url{https://www.mongodb.com/docs/manual/}

\bibitem{atlas} MongoDB, Inc.
  \emph{MongoDB Atlas Documentation}.
  \url{https://www.mongodb.com/docs/atlas/}

\bibitem{motor} MongoDB, Inc.
  \emph{Motor: Asynchronous Python Driver for MongoDB}.
  \url{https://motor.readthedocs.io/}

\bibitem{fastapi} S. Ramírez.
  \emph{FastAPI Documentation}.
  \url{https://fastapi.tiangolo.com/}

\bibitem{pydantic} Pydantic.
  \emph{Pydantic v2 Documentation}.
  \url{https://docs.pydantic.dev/}

\bibitem{uvicorn} Encode.
  \emph{Uvicorn: An ASGI web server for Python}.
  \url{https://www.uvicorn.org/}

\bibitem{react} Meta Open Source.
  \emph{React Documentation}.
  \url{https://react.dev/}

\bibitem{vite} Evan You y contribuidores.
  \emph{Vite: Next Generation Frontend Tooling}.
  \url{https://vitejs.dev/}

\bibitem{recharts} Recharts Group.
  \emph{Recharts: A composable charting library built on React components}.
  \url{https://recharts.org/}

\bibitem{leaflet} V. Agafonkin.
  \emph{Leaflet: an open-source JavaScript library for interactive maps}.
  \url{https://leafletjs.com/}

\bibitem{pandas} The pandas development team.
  \emph{pandas Documentation}.
  \url{https://pandas.pydata.org/docs/}

\bibitem{pytest} H. Krekel y el equipo de pytest.
  \emph{pytest Documentation}.
  \url{https://docs.pytest.org/}

\bibitem{vitest} The Vitest Team.
  \emph{Vitest: A Vite-native testing framework}.
  \url{https://vitest.dev/}

\bibitem{docs-internas} Proyecto WinSports.
  \emph{Documentación interna del repositorio}:
  \texttt{docs/architecture.md}, \texttt{docs/collections.md},
  \texttt{docs/endpoints.md}, \texttt{docs/frontend.md},
  \texttt{docs/admin.md} y \texttt{docs/scripts.md}.

\end{thebibliography}

% ============================================================
\appendix
\section{Anexo: puesta en marcha}
% ============================================================
\begin{lstlisting}[language=bash,caption={Configuración y ejecución}]
# 1. Dependencias
uv sync
cd frontend && npm install && cd ..

# 2. Variables de entorno (.env en la raiz)
#    MONGO_URI=mongodb+srv://<usuario>:<password>@<cluster>.mongodb.net/
#    MONGO_DB_NAME=winsports

# 3. Cargar datos
uv run scripts/load_csv.py --file data/archivo.csv
uv run scripts/test_pipeline.py          # construir derivadas
uv run scripts/test_pipeline.py --HARD   # borrar y reconstruir

# 4. Levantar el dashboard (dos terminales)
uv run uvicorn api.main:app --reload     # API   -> :8000
cd frontend && npm run dev               # Front -> :5173
\end{lstlisting}

\noindent\textbf{Enlaces locales:}
\begin{itemize}[leftmargin=*]
  \item Frontend: \url{http://localhost:5173}
  \item Documentación de la API (Swagger): \url{http://localhost:8000/docs}
  \item Health check: \url{http://localhost:8000/api/health}
\end{itemize}

% ============================================================
\section{Anexo: estructura del repositorio}
% ============================================================
\begin{lstlisting}[language=bash,caption={Árbol de carpetas (resumen)}]
WinSports/
|- api/
|  |- main.py                # app FastAPI + lifespan + routers
|  |- routes/                # overview, qoe, users, alerts, admin
|  +- schemas/               # modelos Pydantic de respuesta
|- config/
|  |- settings.py            # carga de .env (pydantic-settings)
|  +- constants.py           # colecciones, centroides de paises
|- db/
|  |- connection.py          # connect / disconnect / get_db
|  |- csv_ingest.py          # parseo e insercion de CSV (compartido)
|  |- collections/           # validators + setup de cada coleccion
|  |- indexes/               # definicion de indices
|  |- pipelines/             # construccion de sessions y content_stats
|  +- repositories/          # queries por modulo
|- scripts/
|  |- load_csv.py            # carga de events
|  |- test_pipeline.py       # construye derivadas
|  +- verify_connection.py   # prueba de conexion a Atlas
|- frontend/                 # React + Vite (pages, components, layout)
|- data/                     # CSVs de eventos de muestra
|- docs/                     # documentacion tecnica
+- tests/                    # pytest (db/ y api/)
\end{lstlisting}

\end{document}
